\documentclass[11pt]{article}
\usepackage[margin=1in]{geometry}
\usepackage{amsmath,amssymb}
\usepackage{booktabs}
\usepackage[dvipsnames]{xcolor}
\usepackage{hyperref}
\hypersetup{colorlinks=true,linkcolor=blue!50!black,urlcolor=blue!50!black,citecolor=blue!50!black}

\title{Small Samples, Exact Answers:\\
Genome Duplication and the Evolution of Macroscopic Size\\
in Snowflake Yeast, Re-examined with Exact Inference}
\author{David Lindgreen}
\date{October 9, 2026 \\[2pt] \small Version 0.1.0 \quad Preprint, not peer reviewed}

\emergencystretch=3em
\begin{document}
\maketitle

\begin{abstract}
We re-analyse public source data from the long-term snowflake-yeast multicellularity evolution experiment (Tong et
al., 2025)~\cite{tong2025} to ask whether whole-genome duplication is sufficient for macroscopic size. Engineered
tetraploids have a larger 24-hour cluster radius than diploids in both genetic backgrounds tested (PA: $+7.14\,\mu$m, ratio 1.354;
PM: $+17.01\,\mu$m, ratio 1.657), yet evolved PM populations that were already about tetraploid stayed small
(at most 67.7\,$\mu$m through transfer 600 and 53.2\,$\mu$m at transfer 1{,}000) while all five PA populations exceeded
300\,$\mu$m by transfer 1{,}000. Genome duplication therefore helps but is not sufficient. The original analysis used percentile
bootstrap intervals from four replicate strains per group, which are too narrow at that sample size: exact permutation-inversion
intervals are 1.5 to 1.9 times wider (PA 6.46 to 7.92\,$\mu$m; PM 15.65 to 19.21\,$\mu$m), though no conclusion changes. The
preregistered gate passes only because the test was one-sided: the smallest attainable two-sided probability with four against
four replicates is 0.0286, and its Holm adjustment across two backgrounds is 0.0571. An exact five-against-five test shows the PA/PM
divergence is complete by transfer 600 (probability 0.0079, the minimum attainable) but not yet at transfers 200 or 400.
\end{abstract}

\section{Question}
Tong et al.\ evolved snowflake yeast for about 1{,}000 daily transfers and linked genome duplication, cell elongation and packing to
macroscopic multicellularity. The question here is narrow: given the authors' public tables, (i) does engineered tetraploidy
increase cluster size, and (ii) is tetraploidy sufficient for the later macroscopic phenotype? This is a computational reanalysis of
published data, not an independent wet-lab replication.

\section{Data and frozen analysis}
Inputs were derived from the authors' public repository at a pinned commit, each file verified by SHA-256 before use (and re-derived
identically for this paper). The frozen analysis aggregates measurements within biological replicate (four replicate strains per ploidy
per background) before inference, uses an exact one-sided permutation probability over all $\binom{8}{4}=70$ allocations, a Holm
adjustment across the two backgrounds, and a 20{,}000-resample percentile bootstrap of replicate means for the interval. The
longitudinal comparison uses five PA (anaerobic) and five PM (mixotrophic) evolved populations at transfers 200, 400, 600 and 1{,}000.

\section{Results of the frozen analysis}
\begin{table}[h]
\centering\small
\begin{tabular}{lcccc}
\toprule
Background & 2N radius ($\mu$m) & 4N radius ($\mu$m) & Difference ($\mu$m) & Ratio\\
\midrule
PA & 20.157 & 27.299 & +7.142 & 1.354\\
PM & 25.890 & 42.897 & +17.007 & 1.657\\
\bottomrule
\end{tabular}
\caption{Engineered 2N versus 4N 24-hour biomass-weighted mean cluster radius (four replicate strains per group).}
\end{table}
The frozen one-sided exact probability is $1/70=0.0143$ in each background, 0.0286 after Holm adjustment. Evolved lines: all 40
population-time G1 peaks are at least 3.489N, so every sampled population was already about tetraploid; the PA/PM mean-radius ratio is
6.669 at transfer 600 and 8.859 at transfer 1{,}000.

\section{Exact small-sample audit (post-hoc)}
\paragraph{Intervals.}
We replaced the bootstrap interval by the exact permutation-inversion interval: the set of shifts not rejected by the two-sided exact
test~\cite{good2005}. With only four replicates per group a percentile bootstrap undercovers.
\begin{table}[h]
\centering\small
\begin{tabular}{lcc}
\toprule
Background & Frozen bootstrap 95\% ($\mu$m) & Exact inversion 95\% ($\mu$m)\\
\midrule
PA & 6.751 to 7.538 & 6.455 to 7.923\\
PM & 15.976 to 18.318 & 15.649 to 19.210\\
\bottomrule
\end{tabular}
\caption{Interval for the 4N minus 2N radius difference. Both exclude zero.}
\end{table}
\paragraph{The gate depends on one-sidedness.}
With 4 against 4 there are 70 allocations, so the smallest attainable two-sided probability is $2/70=0.0286$, attained here in both
backgrounds. Holm's adjustment across two tests gives $0.0571>0.05$. The preregistered directional hypothesis justifies the one-sided
test (Holm 0.0286), but the conclusion is at the limit of what four replicates per group can establish.
\paragraph{Evolved lines.}
\begin{table}[h]
\centering\small
\begin{tabular}{rcccc}
\toprule
Transfer & Smallest PA ($\mu$m) & Largest PM ($\mu$m) & Separated? & Exact two-sided $p$\\
\midrule
200 & 39.8 & 44.7 & no & 0.325\\
400 & 44.1 & 61.2 & no & 0.095\\
600 & 95.7 & 67.7 & yes & 0.0079\\
1000 & 319.6 & 53.2 & yes & 0.0079\\
\bottomrule
\end{tabular}
\caption{Exact permutation test, five PA against five PM populations; the minimum attainable two-sided probability is $2/252=0.0079$.}
\end{table}

\section{Discussion}
The two lines of evidence answer different questions. The engineered contrast shows an immediate size benefit of duplication in these
strains; the evolved comparison shows that populations which had duplicated their genomes (all G1 peaks at least 3.489N) did not all
become macroscopic, so duplication is not sufficient and other evolved changes separate PA from PM. The divergence emerges between
transfers 400 and 600. The exact tests do not change the direction of any frozen conclusion; they correct the stated precision of
the replicate-level intervals, and they show how thin the formal evidence is at four replicates per group. Neither analysis tests
necessity, a universal threshold, or the authors' mechanical account of entanglement, which the public tables do not measure.

\section{Limitations}
Source-author data from one experimental system; four replicate strains per group; the exact audit covers the primary radius outcome
only (the frozen results also report cell volume and aspect ratio, whose bootstrap intervals have the same limitation); permutation
inversion is evaluated on a fine grid. This is not an independent replication and is not evidence for a general law.

\section*{Reproducibility and AI assistance}
Code, frozen results, the exact audit and tests:\\
\url{https://github.com/lindgreendavid/snowflake-evolution-lab}\\
Interactive laboratory:\\
\url{https://lindgreendavid.github.io/snowflake-evolution-lab/}\\
The analysis code, tests and drafts of this text were produced with substantial assistance from an AI system (Claude, Anthropic);
the author is responsible for the content.

\begin{thebibliography}{9}
\bibitem{tong2025} K.~Tong et al., Genome duplication in a long-term multicellularity evolution experiment,
\emph{Nature} (2025), doi:10.1038/s41586-025-08689-6.
\bibitem{good2005} P.~I. Good, \emph{Permutation, Parametric, and Bootstrap Tests of Hypotheses}, 3rd ed., Springer (2005).
\end{thebibliography}
\end{document}
